\scanpage{61}
\deleted{
% E: text|upper left continuation from previous sheet|d'où un côté privilégié
\deleted{\unsure{deux choix}} d'où un côté privilégié $a$ de $D$ (celui qui rencontre $D_0$), dont je choisis une extrémité $s$, de façon à avoir un repère.
% E: text|left upper middle specialization and half-turn|Je spécialise D en D prime
Je spécialise $D$ en $D'$ vers le \deleted{\unsure{F}} \inserted{sommet} $s'$ le plus proche de $s$, je fais un demi-tour par $s$, et je reviens à $D$, --- un nb pair d'opérations, qui change l'orientation de $D$, \uline{et transforme $s$ en $\unsure{s|s'}$} --- donc c'est la symétrie par \unsure{rapport à} $s$. Si on fait
}
% E: text|left lower new date and definition|4.1 Pour tout chemin
4.1. Pour tout \deleted{\unsure{bout}} chemin \inserted{$c$} combinatoire sur $\Gamma$ --- i.e. chemin formé avec des faces $\overline F_D$ de $(X,K)$, reliant la face $F=\overline F_D$ à une $F'=\overline F_{D'}$, on trouve un iso. de ``transport parallèle''
% E: diagram|left lower|Pol(c) Pol(F) dash Pol(F prime)
\diagram{\begin{tikzcd}
\operatorname{Pol}(c):\quad\operatorname{Pol}(F)\arrow[r,dash] & \operatorname{Pol}(F')
\end{tikzcd}}
% E: text|left bottom intrinsic geometry paragraph|qui du pt de vue
qui du pt de vue de la géométrie intrinsèque de $(X,K,S_K)$ s'interprète par composition successive de ``réflexions'' le long des côtés de réflexion qui interviennent \unsure{au cours sur $c$|au cours de $c$}.
% E: text|right top|L'effet de Pol(c) sur les orientations
L'effet de $\operatorname{Pol}(c)$ sur les orientations
% E: diagram|right upper orientation sets|omegaX(c) omega(F) to omega(F prime)
\diagram{\begin{tikzcd}
\underline\omega_X(c):\quad\underline\omega(F)\simeq\underline\omega(D)\arrow[r] & \underline\omega(F')\simeq\underline\omega(D')
\end{tikzcd}}
% E: text|right upper middle|est aussi l'effet
est aussi l'effet sur l'orientation des \unsure{p.r.|p.v.} par transport parallèle --- donc est indépendant de la parité du nb de pas dans le chemin $c$ ! Si on a un circuit \deleted{$\unsure{\ell}$} $\ell=c$ i.e. $F=F'$ --- on a
% E: diagram|right middle|fundamental group to automorphisms to sign
\diagram{\begin{tikzcd}[column sep=large]
\pi_1(X,F_D)\arrow[r] & \underset{{}\simeq\mathbb D_{2n}}{\operatorname{Aut}(\operatorname{Pol}(D))}\arrow[r,"\substack{\text{effet sur}\\\text{orientation}}"] & \pm1
\end{tikzcd}}
% E: text|right middle introducing character|on désignera par chiX
on désignera par $\chi_X$ le caractère composé
% E: display|right middle character|chiX pi1(X FD) to sign
\[\chi_X:\pi_1(X,F_D)\longrightarrow\pm1\]
% E: text|right middle lower sign interpretation|donc pour un chemin
donc pour un chemin $\ell$, on aura $\chi_{\unsure{X}}(\ell)=+1$ ou $-1$ suivant que par transport parallèle l'orientation de $D$ aura ou non changé. \deleted{\unsure{On}}
% E: text|right lower two other characters|On a deux autres caractères
On a deux autres caractères $\chi_s$ et $\chi_a$ (\unsure{colorations|interversions} sommets et arêtes) en associant à chaque face $F_D$ les deux ensembles à deux éléments
% E: display|right bottom first two-element set|epsilon s(F) epsilon s(D)
\[
\begin{aligned}
\epsilon_s(\unsure{F|D})=\epsilon_s(D)={}&\text{un des }\inserted{\text{deux}}\ n\text{-gones}\\
&\text{inscrits dans }\unsure{\Pi_0|\pi_0}(D)=\operatorname{Pol}(F),
\end{aligned}
\]
% E: display|right bottom second two-element set|epsilon a(D) epsilon a(F)
\[
\begin{aligned}
\epsilon_{\unsure{a|s}}(\unsure{D|F})=\epsilon_a(F)={}&\text{un des deux }n\text{-gones}\\
&\text{circonscrits dans }\operatorname{Pol}(D)=\operatorname{Pol}(F)
\end{aligned}
\]
\mspage{60}
